% ECONOMICS_PROBLEM: {"id": "D78-62", "title": "Politically feasible durable reforms", "jel_code": "D78", "source_id": "OP-0389"}
% !TeX program = xelatex
\documentclass[11pt,a4paper]{article}
\usepackage[margin=23mm]{geometry}
\usepackage{fontspec}
\setmainfont{Latin Modern Roman}
\usepackage{amsmath,amssymb,mathtools}
\usepackage{xcolor,enumitem,booktabs,longtable,array,xurl,hyperref}
\definecolor{muted}{HTML}{53636C}
\hypersetup{colorlinks=true,urlcolor=blue,linkcolor=blue}
\setlength{\parindent}{0pt}
\setlength{\parskip}{5pt}
\newcommand{\R}{\mathbb R}
\newcommand{\E}{\mathbb E}
\newcommand{\N}{\mathbb N}
\newcommand{\ind}{\mathbf 1}
\newcommand{\Var}{\operatorname{Var}}
\newcommand{\Cov}{\operatorname{Cov}}
\newcommand{\supp}{\operatorname{supp}}
\DeclareMathOperator*{\argmin}{arg\,min}
\DeclareMathOperator*{\argmax}{arg\,max}
\newcommand{\entrymeta}[3]{{\sffamily\small\color{muted}\textbf{\texttt{#1}}\quad|\quad #2\par #3\par}}
\newcommand{\Problemref}[1]{\texttt{#1}}
\newcommand{\JELref}[2]{\texttt{#2} (JEL \texttt{#1})}
\begin{document}
\section*{Politically feasible durable reforms}
\textbf{Problem ID:} \texttt{D78-62}\quad \textbf{JEL:} \texttt{D78}\par
% BEGIN ECONOMICS STATEMENT
\label{problem:OP-0389}
{\sffamily\small\color{muted}Original subject group: Political economy and institutions\par}
\label{definitions:problem:30007671}
\textbf{Audit classification:} Background; proposed extension. The World Development Report 2017 primary overview and Main Messages are verified background. Their inspected policy framework does not itself state a further-work passage for this exact question; the research target and variants are editorial extensions.
\subsection*{Definitions and precise research target}
\textbf{Complete editorial problem.} Fix a polity considering a reform with feasible implementation paths \(p\in\mathcal P\). Each path specifies policy content, sequencing, compensation, enforcement, and commitment institutions. Let actor \(i\)'s payoff be \(U_i(p)\), political weight \(\alpha_i\), and declared social welfare weight \(\omega_i\); the two sets of weights need not coincide. Let \(A(p)\) indicate adoption under an explicitly specified bargaining or voting game, and let \(D_T(p)\) indicate continued implementation at horizon \(T\). Define the full-implementation benchmark gain \(G(p)=\sum_i\omega_i\{U_i(p)-U_i(0)\}\), net of real administrative costs and with compensation transfers counted once. The design task is to characterize paths maximizing expected realized welfare, including adoption failure, partial implementation, reversal, and sunk costs, as defined below, subject to fiscal feasibility and credible commitment. Model how compensation affects support, whether promises can be enforced, and how future coalitions may reverse reforms. Empirical work must identify which political constraints bind and test model predictions with relevant reform variation. An efficient policy on paper need not be adoptable or durable. The governance report explicitly poses failures of adoption and implementation; the objective, strategic game, and durability horizon here are our adapted formal research problem, with no universal theorem of optimal reform asserted.

\textbf{Definitions and admissible scope.} Let \(\mathcal P\) be a finite or compact menu of fully specified reform paths. A political model supplies actors, information, timing, voting or bargaining rules, payoffs, and an equilibrium selection rule. Political weights \(\alpha_i\) enter that rule only through an explicit institutional mapping. Social weights \(\omega_i\ge0\) are fixed and normalized separately. Potential payoffs include taxes, compensation, and real implementation costs under a consistent resource budget. If failed or reversed reforms incur costs, \(A(p)D_T(p)G(p)\) is not total welfare. Define \(W(p)=E[\sum_i\omega_i(U_i^{\rm realized}(p)-U_i(0))]\), where realized outcomes include adoption failure, partial implementation, and reversal. Maximize \(W(p)\) subject to budget and sequentially credible commitments; identify the primitives and constraints rather than calling optimization of a known finite menu an unsolved theorem. The target is \(\arg\max_{p\in\mathcal P_{\rm feasible}}W(p)\), where feasibility comprises the stated resource budget and sequential commitment restrictions. For a finite menu, require all expected payoffs finite and a nonempty feasible set. For a compact continuous-path menu, additionally impose upper semicontinuity of \(W\) and closed feasibility to guarantee attainment; otherwise request the supremum and near-optimal paths. Adoption and survival indicators are outcomes of the selected political model, not independent fixed probabilities. The introductory \(G(p)\) describes the fully implemented benchmark; \(U_i^{\rm realized}\) governs actual objective values in every failure, partial-implementation, and reversal state. Identify value bounds when political primitives are uncertain rather than claiming elementary optimization of known inputs is unsolved.
\subsection*{Variants and assumption changes}
\begin{enumerate}
\item Enforceable compensation (editorial specialization): restrict transfers to an enforceable escrow or legal commitment and compare their cost and coalition support with uncompensated reform under the same political rule.
\item Limited commitment and reversal (source-motivated extension): allow future governments to revise transfers and policies after shocks or elections. Compare gradual sequencing and citizen coalition formation using full realized welfare, including costs incurred before reversal.
\end{enumerate}
\subsection*{References and source audit}
\textbf{Support and access.} Primary report page checked. Adoption, implementation and commitment motivate the agenda; the welfare objective and selected strategic game are editorial. \textbf{Locator:} WDR 2017 About, opening questions; Main Messages. \emph{Provenance limit.} The inspected overview presents the report's questions, policy framework, and conclusions. It does not establish the precise editorial problem here as an explicitly published research agenda or open conjecture.
\begin{enumerate}\raggedright
\item\hypertarget{definitions:source:30007671:1}{} World Bank. 2017. \emph{World Development Report 2017: Governance and the Law}. About, opening three paragraphs; Main Messages, first three bullets.
\url{https://www.worldbank.org/en/publication/wdr2017}.
\end{enumerate}

\subsection*{Common conventions: Source mathematical conventions}

These conventions apply to each entry unless explicitly replaced.

\textbf{Models and identification.} A primitive model \(m\) specifies all probability laws, feasible choices, information, equilibrium and measurement rules used by its target. The admissible class \(\mathcal M\) is fixed independently of outcomes. If \(P_O(m)\) is its observable probability law and \(\tau(m)\) the target, its identified set at observable law \(P\) is
\[
 \mathcal I_\tau(P)=\{\tau(m):m\in\mathcal M,\ P_O(m)=P\}.
\]
Point identification means that this set is a singleton. An empty set signals incompatibility of data and assumptions. Coordinate bounds need not describe the joint set. Bounds are sharp only if every retained target is attainable or arbitrarily approachable by compatible models. Finite bounds require explicit restrictions on unobserved quantities; unrestricted confounding, hidden wealth, extrapolation or arbitrary equilibrium selection can yield uninformative sets.

\textbf{Causal interventions.} A potential outcome \(Y_i(p)\) is the outcome under a complete policy regime \(p\), including timing, financing, information and market spillovers. The target population and units are fixed before treatment. With interference, use the complete assignment vector or a justified exposure mapping; individual no-interference notation alone cannot identify an economy-wide intervention. A difference of mean potential outcomes does not require the cross-world joint distribution. Conditioning on post-treatment migration, survival, enrollment or employment changes the estimand and needs separate justification.

\textbf{Statistical statements.} An estimator, confidence set, forecast or test specifies a sampling law, model class and loss/coverage criterion. Uniform \((1-\alpha)\)-coverage means \(\inf_{m\in\mathcal M}\Pr_m\{\tau(m)\in C_n\}\geq1-\alpha\), or an explicitly asymptotic version. The whole parameter space trivially has coverage, so informativeness requires a stated width, risk or power objective. A forecasting improvement is relative to a named benchmark, information set and evaluation population.

\textbf{Equilibrium and optimization.} An equilibrium comprises feasible plans, optimality, beliefs where needed, budget constraints and all market/resource clearing. The primitives or a stated benchmark define each correspondence \(\mathcal E(p;m)\). Nonemptiness is assumed or proved, never inferred just from compact policy sets. A selected-equilibrium target uses a declared selection rule; a robust target quantifies over all permitted equilibria. Use suprema or infima unless attainment is established. An \(\varepsilon\)-optimal policy is feasible and within \(\varepsilon\) of the finite optimum value. Report an empty feasible set separately.

\textbf{Welfare and accounting.} Normative utility, interpersonal normalization, social weights, numeraire, discounting and the use of public receipts are primitives. Behavioral utility need not coincide with the welfare criterion. Household welfare includes ownership income and rents once; adding profits again to compensating variation generally double-counts them. A monetary equivalent is defined by an explicit transfer and price convention, with monotonicity and range conditions. Average effects, mechanism contributions, transfers and welfare losses are different targets. Infinite sums require convergence and dynamic feasible plans require terminal or no-Ponzi conditions.

\textbf{Source versions.} A working-paper date, revision date and journal publication year are distinguished. A DOI for a journal version is not automatically the DOI of an earlier manuscript. Locators refer to the stated version and printed pages where specified. A metadata-only check does not verify a claim about a full-text conclusion. Every entry's source subsection narrows the provenance claim accordingly.
% END ECONOMICS STATEMENT
\end{document}
